
The drift kinetic equation may be solved by different numerical techniques.
Here one of the most natural approach is described, namely the finite
difference method, which is based on the transformation of a differential
equation into an algebraic one. The hyper-volume $\mathcal{V}$ on which the
drift kinetic equation is solved is characterized by three dimensions $
\left( p,\xi ,\psi \right) ,$ and, the time evolution is introduced with the
parameter $t$.

Since the drift kinetic equation is expressed in conservative form, it is
natural to define the grids for the flux $S$ grid as the starting point of
the numerical calculations.
\begin{equation}
S\left( p_{i},\xi _{j},\psi _{l}\right) \rightarrow \left\{ 
\begin{array}{l}
p_{i}\in \left[ 0,p_{\max }\right] ,i\rightarrow \left\{ 0,n_{p}\right\}  \\ 
\xi _{j}\in \left[ -1,1\right] ,j\rightarrow \left\{ 0,n_{\xi }\right\}  \\ 
\psi _{l}\in \left[ 0,\psi _{a}\right] ,l\rightarrow \left\{ 0,n_{\psi
}\right\} 
\end{array}
\right. 
\end{equation}

This grid system defines a set of cells of non-uniform sizescharacterized by
index numbers which, by definition, are integers. The distribution function $%
f$ is represented by its values inside these cells, as shown in Fig. *****.
With this form, the conservative nature of the problem is preserved
numerically, avoiding the occurence of spurious numerical fluxes which lead
to wrong estimate of the solution. 

\begin{equation}
f\left( p_{i+1/2},\xi _{j+1/2},\psi _{l+1/2}\right) \rightarrow \left\{ 
\begin{array}{l}
p_{i+1/2}\in \left[ \frac{\Delta p_{1/2}}{2},p_{\max }-\frac{\Delta %
p_{n_{p}-1/2}}{2}\right] ,i\rightarrow \left\{ 0,n_{p}-1\right\}  \\ 
\xi _{j+1/2}\in \left[ -1+\frac{\Delta \xi _{1/2}}{2},1-\frac{\Delta \xi
_{n_{\xi }-1/2}}{2}\right] ,j\rightarrow \left\{ 0,n_{\xi }-1\right\}  \\ 
\psi _{l+1/2}\in \left[ 0+\frac{\Delta \psi _{1/2}}{2},a-\frac{\Delta \psi
_{n_{r}-1/2}}{2}\right] ,l\rightarrow \left\{ 0,n_{\psi }-1\right\} 
\end{array}
\right. 
\end{equation}

\subsubsection{Grid definitions in momentum space}

Momentum dynamic in $p$ and pitch-angle $\xi $ being independent, the
position of $p_{i+1/2}$ in between $p_{i}$ and $p_{i+1},$ and of $\xi
_{j+1/2}$ in between $\xi _{j}$ and $\xi _{j+1}$ is an arbitrary choice. The
most simple and natural one is to place the point where the distribution
function $f$ is determined exactly in the middle of the cell, 
\begin{equation}
\left\{ 
\begin{array}{l}
p_{i+1/2}=\frac{p_{i}+p_{i+1}}{2} \\ 
\xi _{j+1/2}=\frac{\xi _{j}+\xi _{j+1}}{2}
\end{array}
\right. 
\end{equation}

Then, intervals on the different grids are

\begin{equation}
\left\{ 
\begin{array}{l}
\Delta p_{i+1}=p_{i+3/2}-p_{i+1/2} \\ 
\Delta p_{i+1/2}=p_{i+1}-p_{i} \\ 
\Delta p_{i}=p_{i+1/2}-p_{i-1/2} \\ 
\Delta \xi _{j+1}=\xi _{j+3/2}-\xi _{j+1/2} \\ 
\Delta \xi _{j+1/2}=\xi _{j+1}-\xi _{j} \\ 
\Delta \xi _{j}=\xi _{j+1/2}-\xi _{j-1/2}
\end{array}
\right. 
\end{equation}

where
\begin{eqnarray}
\Delta p_{i+1} &=&p_{i+3/2}-p_{i+1/2} \\
&=&p_{i+3/2}-p_{i+1}+p_{i+1}-p_{i+1/2}  \nonumber \\
&=&\frac{\Delta p_{i+3/2}}{2}+\frac{\Delta p_{i+1/2}}{2}  \nonumber
\end{eqnarray}
\begin{eqnarray}
\Delta \xi _{j+1} &=&\xi _{j+3/2}-\xi _{j+1/2} \\
&=&\xi _{j+3/2}+\xi _{j+1}-\xi _{j+1}-\xi _{j+1/2}  \nonumber \\
&=&\frac{\Delta \xi _{j+3/2}}{2}+\frac{\Delta \xi _{j+1/2}}{2}  \nonumber
\end{eqnarray}

It is important to note that boundary conditions play a crucial role in the
type of equation to be solved, since it corresponds to a Neuman type
problem, where gradients are specified along the perpendicular direction of
the volume of integration. Therefore, $\left( \Delta p,\Delta \xi \right) $
must be specified accurately at each volume boundary. Since $p_{\bot }=0$
corresponds to the cylindrical axis of symmetry of the problem in the
momentum space, 
\begin{equation}
\Delta \xi _{0}=\frac{\Delta \xi _{1/2}+\Delta \xi _{-1/2}}{2}=\Delta \xi
_{1/2}=\xi _{1}-\xi _{0}=2\xi _{1/2}+1
\end{equation}
as $\Delta \xi _{1/2}=\Delta \xi _{-1/2},$ $\xi _{1/2}=\frac{\xi _{1}+\xi
_{0}}{2}$et $\xi _{0}=-1.$ Much in the same way, 
\begin{equation}
\Delta \xi _{n_{\xi }}=\frac{\Delta \xi _{n_{\xi }+1/2}+\Delta \xi _{n_{\xi
}-1/2}}{2}=\Delta \xi _{n_{\xi }-1/2}=\xi _{n_{\xi }}-\xi _{n_{\xi
}-1}=1-\xi _{n_{\xi }-1}
\end{equation}
since $\Delta \xi _{n_{\xi }+1/2}=\Delta \xi _{n_{\xi }-1/2}.$ Furthermore,
since by definition, 
\begin{equation}
\xi _{n_{\xi }-1/2}=\frac{\xi _{n_{\xi }}+\xi _{n_{\xi }-1}}{2}=\frac{1+\xi
_{n_{\xi }-1}}{2}
\end{equation}
one obtains finally

\begin{equation}
\Delta \xi _{n_{\xi }}=2\left( 1-\xi _{n_{\xi }-1/2}\right) 
\end{equation}

A similar technique may be used for the momentum grid $p$. Hence 
\begin{equation}
\Delta p_{0}=\frac{1}{2}\left( \Delta p_{1/2}+\Delta p_{-1/2}\right) =\Delta %
p_{1/2}
\end{equation}
using $\Delta p_{-1/2}=\Delta p_{1/2}.$ Since $\Delta p_{1/2}=p_{1}-p_{0}$,
and $p_{1/2}=\frac{p_{1}+p_{0}}{2}$ one obtains finally,
\begin{equation}
\Delta p_{0}=2p_{1/2}
\end{equation}
because $p_{0}=0.$ Much in the same way,
\begin{equation}
\Delta p_{n_{p}}=\Delta p_{\max }=\frac{1}{2}\left( \Delta p_{n_{p}+1/2}+%
\Delta p_{n_{p}-1/2}\right) =\Delta p_{n_{p}-1/2}=p_{n_{p}}-p_{n_{p}-1}
\end{equation}
and since by definition $p_{n_{p}+1/2}=\frac{p_{n_{p}}+p_{n_{p}-1}}{2}=\frac{%
p_{\max }+p_{n_{p}-1}}{2},$ it results 
\begin{equation}
\Delta p_{n_{p}}=2\left( p_{n_{p}}-p_{n_{p}-1/2}\right) 
\end{equation}
or in equivalent manner, 
\begin{equation}
\Delta p_{\max }=2\left( p_{\max }-p_{n_{p}-1/2}\right) 
\end{equation}

Note that the exact value of $\Delta p_{\max }$ has no importance, since the
numerical problem simplifies into an hyperbolic one at this particular
point, and therefore, it is not necessary to specify boundary conditions.

\subsubsection{Grid definitions in configuration space}

By an appropriate bounce averaging, the number of dimension is reduced to
one, i.e. the radial one. As for the dynamics in momentum space, the point
where the distribution function $f$ is determined is chosen exactly in the
middle of the cell, 

\begin{equation}
\psi _{l+1/2}=\frac{\psi _{l}+\psi _{l+1}}{2}
\end{equation}
and the intervals on the different grids are,
\begin{equation}
\left\{ 
\begin{array}{l}
\Delta \psi _{l+1}=\psi _{l+3/2}-\psi _{l+1/2} \\ 
\Delta \psi _{l+1/2}=\psi _{l+1}-\psi _{l} \\ 
\Delta \psi _{l}=\psi _{l+1/2}-\psi _{l-1/2}
\end{array}
\right. 
\end{equation}
where
\begin{eqnarray}
\Delta \psi _{l+1} &=&\psi _{l+3/2}-\psi _{l+1/2} \\
&=&\psi _{l+3/2}-\psi _{l+1}+\psi _{l+1}-\psi _{l+1/2}  \nonumber \\
&=&\frac{\Delta \psi _{l+3/2}}{2}+\frac{\Delta \psi _{l+1/2}}{2}  \nonumber
\end{eqnarray}

Much in the same way,  $\Delta \psi $ must be specified at the boundaries of
the domain of integration. Using a similar approach as for the momentum  $p$%
, one finds that 
\begin{equation}
\Delta \psi _{0}=2\psi _{1/2}
\end{equation}
since $\psi _{0}=0$, and 
\begin{equation}
\Delta \psi _{n_{\psi }}=2\left( \psi _{n_{\psi }}-\psi _{n_{\psi
}-1/2}\right) =2\left( \psi _{a}-\psi _{n_{\psi }-1/2}\right) 
\end{equation}
en remarquant que $\psi _{n_{\psi }}=\psi _{a}$. 

\subsubsection{Link between momentum and configuration grids}

When the bounce averaged drift kinetic equation is solved, the distribution
function $f^{\left( 0\right) }\equiv f\left( B=B_{\min }\right) $ is
determined by definition at the poloidal angle value where the magnetic
field is minimal on a given magnetic flux surface. It is indeed the only
common poloidal point on the magnetic flux surface to all particle
trajectories, whatever they are trapped or or not. In momentum space, the
dynamics of trapped and passing particles is characterized by an interval
boundary which is defined by the relation
\begin{equation}
\xi _{0T}^{2}\left( \psi \right) =1-\frac{B_{\min }\left( \psi \right) }{%
B_{\max }\left( \psi \right) }  \label{Trapcond}
\end{equation}

From relation (\ref{Trapcond}), it is clear that momentum and radial grids
are no more fully independent, if one enforces for ensuring numerical
accuracy, that at each radial position $\psi _{l}$ corresponds an index
number $j_{l}$ of the picth-angle grid, so that 
\begin{equation}
\xi _{0T,j_{l}}^{2}\left( \psi _{l}\right) =1-\frac{B_{\min }\left( \psi
_{l}\right) }{B_{\max }\left( \psi _{l}\right) }
\end{equation}

Hence, if the number of radial point $n_{\psi }$ is arbitrary, the number of
points for the pitch-angle grid has a lower limit 
\[
\xi _{0T}^{2}=\frac{2\varepsilon }{1+\varepsilon }
\]
$\varepsilon =r/R_{0}$ etant l'inverse du rapport d'aspect. De cette
relation il est clair que la grille angulaire en $\xi $ et la grille radiale
en $r$ ne sont plus independantes si l'on impose qu'a chaque point de la
grille radial

e en flux doit correspondre un point de la grille angulaire en flux
caracterisant pour ce lieu la frontiere entre les domaines circulant et
pieges. La relation entre les points de grille vaut 
\[
\left| \xi _{0T,j_{l}}\right| =\sqrt{\frac{2\varepsilon _{l}}{1+\varepsilon
_{l}}}
\]
ou $\varepsilon _{l}=r_{l}/R_{0}$. De cette maniere, on assure une coherence
entre les grilles en flux dans les differents espaces consideres, une
condition necessaire pour un calcul numerique de grande precision. Pour une
grille radiale en flux donnee, il convient donc d'adapter la grille
angulaire en flux de telle sorte que cette condition de frontiere entre les
particules circulantes et piegees soit remplie a chaque position radiale.
Par principe donc, la grille radiale des flux est arbitraire et l'espacement
des points de grille ne depend que de la precision que l'on cherche
localement en $r$ en fonction des gradients du probleme etudie. Si cette
grille est caracterisee par $n_{r}+1$ points, la grille angulaire en flux
aura donc au moins $n_{\xi }>6n_{r}+7$, chaque frontiere etant caracterisee
par 3 points de grille, un place exactement sur cette frontiere, et les deux
autres places symetriquement de part et d'autre du point central. Ce nombre
est multiplie par deux pour tenir compte des domaines $\xi >0$ et $\xi \leq 
0,$ les points de grille supplementaires correspondant respectivement a $%
\left| \xi \right| =1$ (4 points de grille) et $\xi =0$ (3 points de grille)
pour assurer une bonne precision des calculs aux extremites du domaine
angulaire. Entre ces valeurs imposees par la moyenne sur la periode de
rebond, il est loisible de choisir n'importe quelle repartition des points
de grille angulaire en flux, selon la precision recherchee pour les calculs.
Par exemple, pour 21 points radiaux en flux, le nombre minimal de points sur
la grille angulaire en flux doit donc etre superieur a 133.
 